%% 第三章 数据集构建与数据工程
\chapter{数据集构建与数据工程}

数据是目标检测系统的基石。本章介绍本文使用的三类数据：公开的 Det-Fly 无人机检测数据集、自建的 AprilTag 标定与定位评测数据集、以及仿真域自动标注数据集。重点说明数据的来源、格式转换、划分策略、统计特性与评测口径，并讨论数据泄漏、类别不平衡与尺度分布等容易被忽视的问题。

\section{Det-Fly 数据集}

\subsection{数据集来源与内容}

Det-Fly\cite{detfly} 是反无人机视觉检测领域常用的公开数据集，由多段无人机飞行视频帧组成，标注格式为 YOLO 文本格式，类别为单一类别 UAV。数据覆盖多种背景（天空、树林、建筑、山地）、多种机型与多种飞行状态（悬停、平飞、盘旋、远近变化）。图像分辨率与拍摄设备在不同序列间存在差异，整体以中远距离拍摄为主，目标在画面中的占比普遍较小。

与 COCO 等通用数据集相比，Det-Fly 的显著特点是\textbf{目标尺度分布极度偏向小目标}。本文对转换后的 13271 个标注框做了统计，结果如表~\ref{tab:detfly_stats} 所示。

\begin{table}[htbp]
\centering
\caption{Det-Fly 数据集标注统计}
\label{tab:detfly_stats}
\small
\begin{tabular}{lcc}
\toprule
\textbf{统计项} & \textbf{数值} & \textbf{占比} \\
\midrule
标注文件数 & 13271 & -- \\
标注框总数 & 13270 & 100\% \\
类别数 & 1（UAV） & 100\% \\
\midrule
small（$<32^2$\,px） & 1490 & 11.2\% \\
medium（$32^2\sim96^2$\,px） & 8067 & 60.8\% \\
large（$>96^2$\,px） & 3713 & 28.0\% \\
\midrule
目标框面积最小值 & 0.0013\% 画面面积 & -- \\
目标框面积中位数 & 0.0554\% 画面面积 & -- \\
目标框面积 $<1\%$ 画面 & -- & 99.1\% \\
目标框面积 $<5\%$ 画面 & -- & 100.0\% \\
\bottomrule
\end{tabular}
\end{table}

由表可见，虽然按 COCO 面积口径仅 11.2\% 的框属于 small，但这是相对于各图像自身分辨率计算的结果；若统一折算到 1280$\times$720 的推理分辨率，绝大多数目标在画面中的面积不足 1\%。这解释了为什么通用检测模型在该数据集上的 AP\_small 极低，也说明了增加高分辨率检测层的必要性。

\subsection{格式转换}

原始数据以视频帧与独立标注文件的形式组织，本文编写了转换脚本将其整理为 YOLO 训练所需的目录结构：

\begin{verbatim}
dataset/
  images/
    train/  val/
  labels/
    train/  val/
  data.yaml
\end{verbatim}

每张图像对应一个同名 \texttt{.txt} 标注文件，每行格式为

\begin{verbatim}
<class_id> <x_center> <y_center> <width> <height>
\end{verbatim}

其中后四项均为相对图像宽高的归一化值，取值范围 $[0,1]$。转换过程中做了三项校验：坐标越界裁剪、零面积框剔除、图像与标注一一对应检查。最终保留 13270 个有效标注框。

\subsection{序列级划分与数据泄漏}

数据划分是影响评测可信度的关键因素。若随机按帧划分，同一段视频的相邻帧可能同时进入训练集与验证集，由于相邻帧在背景、光照、目标姿态上高度相似，模型实际上"见过"验证样本，导致指标虚高。这种数据泄漏在小目标检测中尤为严重，因为模型可能记住背景纹理而非学习目标特征。

本文按\textbf{序列}划分：以序列编号中的 010 与 020 为界，010 序列全部用于训练，020 序列全部用于验证。这样可以保证验证集中的背景、光照与飞行场景在训练集中未出现过，评测结果更接近真实部署时的泛化性能。划分比例约为 8:2，与常用实践一致。

需要说明的是，序列级划分会使验证集指标低于随机划分，这是正常现象。本文在报告中统一注明"序列级划分"，以便读者正确理解指标含义。

\subsection{数据增强}

训练阶段采用 Ultralytics 默认的数据增强策略，包括：

（1）\textbf{Mosaic}：将四张图像随机拼接为一张，增加目标在不同位置、不同尺度下的出现概率，同时变相扩大批量大小；

（2）\textbf{HSV 扰动}：随机调整色调、饱和度与明度，提升对光照变化的鲁棒性；

（3）\textbf{随机翻转与平移}：水平翻转概率 0.5，平移比例 0.1；

（4）\textbf{尺度缩放}：随机缩放比例 0.5，模拟目标远近变化。

在仿真域微调阶段，本文\textbf{关闭 Mosaic}。原因是仿真数据的尺度分布是精心分层设计的，Mosaic 拼接会破坏这一分布，使近距离大目标被缩小、远距离小目标被放大，从而削弱微调对死区的针对性。

\section{自建 AprilTag 数据集}

\subsection{采集目的}

公开数据集缺少与本文相机、镜头、安装方式一致的标定数据，难以评估单目定位的绝对精度。为此，本文自建了一套基于 AprilTag 的评测数据集，用于：

（1）标定相机内参，获得 $f_x,f_y,c_x,c_y$ 与畸变系数；

（2）在已知距离下评估测距误差；

（3）在已知角度下评估方位角与俯仰角误差；

（4）验证状态机在目标静止、缓慢移动与遮挡条件下的行为。

\subsection{硬件与采集流程}

采集使用普通 USB 相机，分辨率 1920$\times$1080，水平视场角约 60\textdegree。标定板为 9$\times$6 的棋盘格，方格边长 25\,mm，打印后粘贴在硬质平板上以保证平整。AprilTag 采用 36h11 家族、ID 为 0、边长 16\,cm 的标签，打印后粘贴在无人机机身侧面。

采集流程如下：

（1）\textbf{内参标定}：手持相机从不同角度拍摄棋盘格 30 余张，覆盖画面中心与四角、不同倾斜角度与不同距离，使用 OpenCV 的 \texttt{calibrateCamera} 求解内参与畸变系数，重投影误差控制在 0.5\,px 以内；

（2）\textbf{距离评测}：将 AprilTag 固定在无人机上，用卷尺量取相机到目标的真实距离，分别在 2、3、4、5、6、8、10\,m 处采集静止图像，每个距离采集 20 帧；

（3）\textbf{角度评测}：保持距离不变，将目标沿水平方向与垂直方向移动到画面不同位置，记录真实角度（由几何关系计算），每个角度采集 10 帧；

（4）\textbf{动态评测}：手持目标缓慢移动，模拟真实飞行，采集连续视频用于状态机与滤波测试。

所有采集场景记录在 \texttt{meta.json} 中，包含时间、距离、角度、光照条件与备注，便于后续复核。

\subsection{真值精度分析}

卷尺量距的读数误差约为 $\pm1$\,cm，在 10\,m 处相对误差约 0.1\%，远小于单目测距本身的误差，因此可以作为真值使用。角度真值由目标中心在画面中的几何位置与实测距离反算，误差主要来自目标中心的手工标注，约 $\pm2$\,px，对应约 0.1\textdegree，同样小于算法误差。需要说明的是，AprilTag 的检测中心与无人机几何中心存在小幅偏移，本文在评测中对偏移做了补偿。

\section{仿真域自动标注数据集}

\subsection{构建动机}

如第一章所述，直接用 Det-Fly 训练的模型在仿真画面中会出现严重死区。为通过微调消除死区，需要一批"仿真画面 + 真值标注"的数据。人工标注 450 帧仿真图像成本较高，且框的边界一致性难以保证。本文利用仿真环境的可控性，设计了零人工标注的自动构建流程。

\subsection{自动标注原理}

jMAVSim 的护航相机通过 \texttt{setViewerTargetObject} 接口始终注视目标机，因此目标中心在画面中的位置近似等于图像中心：

\begin{equation}
c_x \approx c_{x0},\qquad c_y \approx c_{y0}+\Delta_y,
\label{eq:autolabel_center}
\end{equation}

其中 $c_{x0},c_{y0}$ 为图像中心，$\Delta_y$ 为标定得到的纵向偏移（本文实测约 $-17$\,px，源于相机安装角与投影中心的偏差）。

目标框的高度由真值距离与等效尺寸反算：

\begin{equation}
h_{px} = \frac{f\cdot D}{Z_{true}},
\label{eq:autolabel_h}
\end{equation}

其中 $f$ 为以像素为单位的焦距，$D$ 为无人机等效尺寸，$Z_{true}$ 为真值距离。框宽按目标长宽比 $r_w$ 放大：

\begin{equation}
w_{px} = r_w\cdot h_{px},\qquad r_w \approx 2.62.
\label{eq:autolabel_w}
\end{equation}

为避免框过紧导致训练时目标被裁切，本文对框做了 1.15 倍的尺度放大，以覆盖旋翼与起落架。最终标注框为

\begin{equation}
\left(c_x-\frac{w_{px}}{2},\,c_y-\frac{h_{px}}{2},\,c_x+\frac{w_{px}}{2},\,c_y+\frac{h_{px}}{2}\right).
\label{eq:autolabel_box}
\end{equation}

\subsection{距离分层采样}

死区的两端分别是近距离大尺度与远距离小尺度，因此采样必须覆盖完整距离范围。本文按表~\ref{tab:sim_sampling} 进行分层采样。

\begin{table}[htbp]
\centering
\caption{仿真微调数据距离分层采样}
\label{tab:sim_sampling}
\small
\begin{tabular}{lccc}
\toprule
\textbf{距离区间(m)} & \textbf{帧数} & \textbf{采样目的} & \textbf{对应问题} \\
\midrule
3--6 & 150 & 近距离大尺度 & 训练分布外（OOD） \\
6--7 & 80 & 过渡区 & 死区边缘 \\
7--9 & 100 & 原有效窗口 & 保持既有能力 \\
9--12 & 120 & 远距离小尺度 & 低于初锁门限 \\
\midrule
合计 & 450 & -- & -- \\
\bottomrule
\end{tabular}
\end{table}

采集时由脚本控制目标机按预设航线飞行，同时记录 TCP 帧与 UDP 真值，按上述规则实时生成标注。训练集与验证集按 360:90 划分，划分时保证同一段连续飞行不跨集，避免泄漏。

\subsection{数据质量核验}

自动标注的质量取决于三项假设：目标中心在光轴附近、等效尺寸准确、真值时间戳与帧对齐。本文对每项假设做了核验。

（1）\textbf{中心假设}：抽取 50 帧人工检查，目标中心与标注中心的水平偏差中位数为 3\,px，垂直偏差中位数为 5\,px，均在框高的一半以内，说明假设成立。

（2）\textbf{尺寸假设}：等效尺寸由真值距离与实测框高反算得到，本文在 3--10\,m 范围内做了 20 组测量，反算尺寸的变异系数小于 5\%，说明尺寸假设稳定。

（3）\textbf{时间对齐}：真值广播频率为 50\,Hz，帧率为 30\,fps，同一帧对应的时间戳误差小于 20\,ms，在目标速度 2\,m/s 时对应位移小于 4\,cm，对框位置的影响可以忽略。

\section{评测口径与可复现性}

为避免不同实验之间的指标不可比，本文统一了以下口径。

（1）\textbf{检测指标}：mAP50、mAP50-95、Precision、Recall 采用 Ultralytics 官方验证协议，取验证集上 mAP50-95 最优的 epoch 报告；AP\_small/medium/large 采用 COCO 面积口径、IoU=0.5、101 点插值，由独立脚本计算。

（2）\textbf{速度指标}：FPS 为 RTX 4060、batch=1、帧预载内存、包含预处理与后处理的口径；不同模型的测量脚本与输入尺寸完全一致。

（3）\textbf{闭环指标}：检出率为检出帧数除以全部处理帧数（含搜索帧），而非仅统计锁定后的帧；测距 RMSE 在检出帧上计算；稳态口径为连续检出不少于 5 帧后的样本，以体现滤波收敛后的性能。

（4）\textbf{角度指标}：采用视轴偏角口径，即目标相对相机光轴的角度误差，而非世界坐标系下的绝对方位。这一区分在仿真中尤为重要：护航相机会自动注视目标，世界系方位角随相机偏航而变化，若混淆两种口径会得出错误的"19\textdegree 误差"结论。

（5）\textbf{随机性控制}：训练使用固定随机种子；仿真航迹由种子化随机数生成；所有实验重复至少两次，报告中位数或均值。

\section{本章小结}

本章介绍了本文的三类数据及其工程处理。Det-Fly 是检测模型的主训练集，其目标尺度极度偏小，且必须按序列划分以避免数据泄漏；自建 AprilTag 数据集用于相机标定与定位精度评估，提供了可信的绝对真值；仿真自动标注数据集利用护航相机注视特性实现了零人工标注，并通过距离分层采样覆盖死区两端。最后，本章统一了检测、速度、闭环与角度四类指标的口径，为后续章节的实验结果提供了可比基础。
